%==============================================================================
\section{Introduction}
%==============================================================================

Threat modeling is a cornerstone of secure software development~\cite{shostack2014threat}, requiring practitioners to systematically identify threats, analyze attack vectors, and propose mitigations. As cloud-native architectures grow in complexity, manual threat modeling becomes increasingly impractical---a single system may expose hundreds of attack surfaces across multiple services, each requiring domain-specific security expertise.

Large language models (LLMs) offer a path toward automating this process. Recent multi-agent systems can generate threat statements, construct attack trees~\cite{schneier1999attack}, and map findings to frameworks like MITRE ATT\&CK~\cite{strom2018mitre}. However, a fundamental question remains unanswered: \emph{how do we measure whether automated threat models are any good?}

Unlike code generation---where correctness can be verified by compilation and test suites---threat modeling quality is inherently multi-dimensional and subjective. A generated attack tree may be syntactically valid yet unrealistic, structurally complete yet missing critical attack phases, or technically accurate yet unhelpful for defenders. No existing benchmark or evaluation methodology addresses this challenge.

We identify three gaps in the current landscape:

\textbf{No standardized evaluation dimensions.} Existing work evaluates threat modeling tools informally, typically through case studies or user surveys. There is no agreed-upon set of quality dimensions, scoring scales, or evaluation protocols.

\textbf{No separation of capabilities.} Automated threat modeling involves distinct sub-tasks---generating threat statements, constructing attack trees, mapping to TTP frameworks---each requiring different evaluation criteria. Existing approaches conflate these into a single quality judgment.

\textbf{No infrastructure for iterative improvement.} Without structured data collection and scoring, there is no path from evaluation to optimization. Improving prompt quality requires ground truth datasets that do not yet exist.

We address these gaps with an evaluation framework for \threatforest{}, a multi-agent threat modeling system. Our contributions are:

\begin{enumerate}
    \item \textbf{Multi-dimensional evaluation taxonomy.} We decompose threat modeling quality into three capabilities---threat statement generation, attack tree generation, and TTP matching---with 12 score dimensions across a standardized 5-point categorical scale.

    \item \textbf{Hybrid evaluation methodology.} We combine automated structural metrics (node count, path coverage, phase detection, syntax validation) with SME scoring protocols, enabling both scalable measurement and expert validation.

    \item \textbf{Tracing infrastructure.} We implement the framework as a production tracing system built on OpenTelemetry and Langfuse, capturing every agent interaction with structured metadata, scores, and export pipelines for dataset curation.

    \item \textbf{Baseline measurements.} We report the first quantitative baselines for automated threat modeling quality across \todo{N} diverse domains.

    \item \textbf{Optimization roadmap.} We describe how the evaluation framework enables iterative prompt optimization via DSPy's MIPROv2~\cite{opsahlong2024miprov2}, where SME-scored traces feed directly into reflective optimization cycles.
\end{enumerate}
